\chapter{Información General del Proyecto}
\label{cap:01_informacion_general}

El presente capítulo proporciona una descripción general del proyecto \textbf{Rhyzome}, incluyendo su identificación, propósito, equipo de trabajo, organización de responsabilidades, acompañamiento académico y duración. Esta información permite contextualizar el desarrollo del proyecto y establecer las condiciones generales bajo las cuales se llevó a cabo.

% ============================================================
% 1.1. NOMBRE DEL PROYECTO
% ============================================================

\section{Nombre del proyecto}

El proyecto se denomina \textbf{Rhyzome --- Generador Procedural de Mapas 2D para Videojuegos}.

Rhyzome es una aplicación web orientada a la generación procedural de escenarios bidimensionales para videojuegos, mediante la integración de algoritmos de generación procedural y técnicas de aprendizaje automático (\textbf{Machine Learning}).

La solución permite generar diferentes tipos de escenarios, como bosques, desiertos, ciudades y mazmorras, utilizando algoritmos como Perlin Noise y Autómatas Celulares. Adicionalmente, incorpora un componente inteligente que analiza las características de los mapas generados y apoya su selección de acuerdo con niveles de complejidad definidos.

El proyecto busca proporcionar una herramienta que facilite la creación y experimentación con escenarios bidimensionales, integrando generación procedural, evaluación de características y técnicas de inteligencia artificial.

% ============================================================
% 1.2. EQUIPO DE TRABAJO
% ============================================================

\section{Equipo de trabajo}

El desarrollo de Rhyzome se realiza como un proyecto académico individual en el marco del programa de maestría de la Universidad del Valle.

El equipo de trabajo está conformado por \textbf{Christian Andres Morales  Mera}, responsable de la planificación, diseño, implementación, validación y documentación del sistema.

La ejecución del proyecto comprende actividades relacionadas con el desarrollo de software, la implementación de algoritmos procedurales, la experimentación con modelos de aprendizaje automático y la gestión de actividades mediante prácticas ágiles.

% ============================================================
% 1.3. INTEGRANTES Y ROLES SCRUM
% ============================================================

\section{Integrantes y roles SCRUM}

Para la organización y seguimiento del desarrollo se adoptan prácticas de la metodología  ágil Scrum, adaptadas a las características de un proyecto académico individual.

Debido a que el proyecto cuenta con un único desarrollador, las responsabilidades asociadas a los roles de Scrum se concentran en una misma persona, sin constituir un equipo Scrum tradicional.

La Tabla~\ref{tab:roles_scrum} presenta la
distribución de responsabilidades.


\begin{table}[htbp]
    \centering

    \caption{Roles y responsabilidades en Rhyzome}
    \label{tab:roles_scrum}

    \renewcommand{\arraystretch}{1.5}

    \begin{tabularx}{\textwidth}{
        |>{\raggedright\arraybackslash}p{3cm}|
         >{\raggedright\arraybackslash}X|
    }
        \hline

        \rowcolor{gray!15}
        \textbf{Rol} &
        \textbf{Responsabilidades} \\

        \hline

        \textbf{Product Owner} & Definir y priorizar los requerimientos, gestionar el Product Backlog y establecer los objetivos de cada entrega. \\

        \hline

        \textbf{Scrum Master} & Organizar las actividades de planificación, seguimiento y retrospectiva, identificando impedimentos y oportunidades de mejora. \\

        \hline

        \textbf{Developer} & Diseñar, implementar, integrar, probar y documentar los componentes del sistema, incluyendo frontend, backend y modelos de aprendizaje automático. \\

        \hline

    \end{tabularx}

    \par\smallskip
    \raggedright

    \textit{Nota.} Las responsabilidades descritas son asumidas por Christian Andrés Morales Mera como adaptación de Scrum a un proyecto individual.

\end{table}


La gestión de actividades se apoya en herramientas como Jira para la planificación y seguimiento de historias de usuario, GitHub para el control de versiones y la integración continua, y Confluence para la organización de documentación y evidencias del proyecto.

% ============================================================
% 1.4. PROFESOR ASESOR
% ============================================================

\section{Profesor asesor}

El acompañamiento académico del proyecto está a cargo de la profesora \textbf{Deisy Chaves Sánchez}, de la Universidad del Valle.

Su participación se enmarca en el proceso de orientación y seguimiento académico del trabajo integrador, contribuyendo a la revisión del avance del proyecto y los entregables establecidos en el contexto de la asignatura.

% ============================================================
% 1.5. DURACIÓN TOTAL DEL PROYECTO
% ============================================================

\section{Duración total del proyecto}

El proyecto Rhyzome se desarrolla de manera progresiva a través de las actividades académicas y técnicas correspondientes al trabajo integrador.

Su ejecución contempla las etapas de planificación, análisis de requerimientos, diseño de arquitectura, implementación de algoritmos procedurales, desarrollo del componente inteligente, integración de la aplicación, pruebas, validación y documentación final.

La gestión del desarrollo se organiza mediante iteraciones y entregas incrementales, permitiendo incorporar mejoras funcionales y técnicas
durante el ciclo de vida del proyecto.

\begin{description}

    \item[Fecha de inicio:]
    Enero de 2026.

    \item[Fecha de finalización:]
    Diciembre de 2026.

    \item[Duración total:]
    Aproximadamente 12 meses de desarrollo.

\end{description}

La planificación y el seguimiento de las iteraciones, los sprints y los releases se presentan en el Capítulo~\ref{cap:05_gestion_agil}.